\begin{helper}[Elementary facts about the gradient field of a smooth compactly supported function and about $L^2$ fields]
Let $n\in\N$.  Throughout, $\operatorname{supp}\Phi$ denotes the \emph{closed}
support of $\Phi$; the gradient field of a function $\Phi$ on $\R^n$ is the
$\R^n$-valued map
$\nabla\Phi\colon x\mapsto\bigl(\partial_1\Phi(x),\dots,\partial_n\Phi(x)\bigr)$,
assembled from the directional derivatives of the Fréchet differential along
the standard basis vectors, so that $|\nabla\Phi(x)|$ is the Euclidean norm of
that vector and not an operator norm; $v_i$ is the $i$-th coordinate of
$v\in\R^n$; and the exponent $2$ is read as an extended-real exponent where an
$L^2$ membership is asserted.  Then all four of the following hold.
\begin{enumerate}
\item For every $\Phi\colon\R^n\to\R$ which is smooth and has compact support,
  the field $\nabla\Phi$ is continuous, has compact support, and satisfies
  $\nabla\Phi(x)=0$ for every $x\notin\operatorname{supp}\Phi$.
\item For every $v\in\R^n$ and every index $i$, $|v_i|\le|v|$.
\item For every measurable $S\subseteq\R^n$ and every
  $v\colon\R^n\to\R^n$ with $v\in L^2(S)$, the function $x\mapsto|v(x)|^2$ is
  integrable on $S$.
\item $1\le2$ as extended-real exponents.
\end{enumerate}
No hypothesis relates the four statements; they are bundled because the
approximation passages that use them use all four.
\end{helper}
\begin{proof}
(1)  Since $\Phi$ is smooth, its differential $x\mapsto d\Phi(x)$ is continuous
(Preamble, continuity of the derivative of a $C^\infty$ map), and evaluation of
a continuous linear map at a fixed vector is continuous in the map, so each
coordinate $x\mapsto\partial_i\Phi(x)=d\Phi(x)[e_i]$ is continuous; assembling
finitely many continuous coordinates into a vector is continuous, whence
$\nabla\Phi$ is continuous.  If $x\notin\operatorname{supp}\Phi$ then
$d\Phi(x)=0$, because the differential of a function vanishing on a
neighbourhood of $x$ vanishes at $x$ and the complement of the closed support
is open; hence every coordinate of $\nabla\Phi(x)$ is $0$ and
$\nabla\Phi(x)=0$.  For the compact support, the set
$\{x:\nabla\Phi(x)\neq0\}$ is contained in $\operatorname{supp}\Phi$ by what
was just proved, so its closure is a closed subset of the compact set
$\operatorname{supp}\Phi$ and is therefore compact.

(2)  Writing $v_i=\langle e_i,v\rangle$ for the standard basis vector $e_i$,
the Cauchy--Schwarz inequality (Preamble) gives
$|v_i|=|\langle e_i,v\rangle|\le|e_i|\,|v|=|v|$, since $|e_i|=1$.

(3)  If $v\in L^2(S)$ then $x\mapsto|v(x)|^{2}$ is integrable on $S$ by the
characterisation of $L^q$ membership through integrability of the $q$-th power
of the norm (Preamble), the exponent being finite and nonzero.  That
characterisation produces the real power $|v|^{2}$, which coincides with the
square, so the asserted integrability follows.

(4)  $1\le2$ holds in $\R$, hence between the corresponding extended-real
exponents.
\end{proof}
